\subsection{Уравнение второго рода с конечномерным оператором}

Рассматривается уравнение

\stepcounter{equation}
\begin{equation}\tag{I}\label{eq:deg-main}
  u(x) - \int\limits_{0}^{\ell} K(x,y)\, u(y)\, dy = f(x) ,
  \qquad\text{или}\qquad
  u - Au = f \ \text{ в } H ,
\end{equation}
причём ядро \textbf{вырождено}:
\[
  K(x,y) = \sum_{i=1}^{m} \varphi_i(x)\, \overline{\psi_i(y)} ,
\]
где системы $\{\varphi_i\}$ и $\{\psi_i\}$ линейно независимы.

\subsubsection*{Оператор конечномерен}

\[
  (Au)(x)
  = \int\limits_{0}^{\ell} \sum_{i=1}^{m} \varphi_i(x)\overline{\psi_i(y)}\,
    u(y)\, dy
  = \sum_{i=1}^{m} \left[\ \int\limits_{0}^{\ell} u(y)\,\overline{\psi_i(y)}\,
    dy \right] \varphi_i(x) ,
\]
то есть

\stepcounter{equation}
\begin{equation}\tag{2}\label{eq:deg-A}
  Au = \sum_{i=1}^{m} (u, \psi_i)\, \varphi_i ,
  \qquad \varphi_i, \psi_i \in H .
\end{equation}
Образ $A$ лежит в линейной оболочке $\varphi_1,\dots,\varphi_m$, значит $A$ —
конечномерный оператор.

\subsubsection*{Сведение к системе}

Подставляя~\eqref{eq:deg-A} в~\eqref{eq:deg-main}:
\[
  u = \sum_{i=1}^{m} \underbrace{(u, \psi_i)}_{=:\ u_i}\, \varphi_i + f .
\]
Домножим это равенство скалярно на $\psi_i$ — получим замкнутую систему
относительно чисел $u_i$:
\[
  u_i = (u, \psi_i)
      = \left(\sum_{j=1}^{m} u_j \varphi_j + f,\ \psi_i\right)
      = \sum_{j=1}^{m} a_{ij}\, u_j + f_i ,
\]
где
\[
  a_{ij} = (\varphi_j, \psi_i) , \qquad f_i = (f, \psi_i) .
\]
В матричной записи

\stepcounter{equation}
\begin{equation}\tag{I$'$}\label{eq:deg-sys}
  \vec u - \mathfrak{A}\, \vec u = \vec f ,
  \qquad \mathfrak{A} = [a_{ij}] \in \CC^{m\times m} .
\end{equation}

\inputfig{figures/fig-correspondence}

\begin{proposition}\label{prop:deg-correspondence}
Решения~\eqref{eq:deg-main} и решения~\eqref{eq:deg-sys} находятся во взаимно
однозначном соответствии:
\[
  u \ \longmapsto\ \vec u = \big((u,\psi_1), \dots, (u,\psi_m)\big) ,
  \qquad
  \vec u \ \longmapsto\ u = \sum_{i=1}^{m} u_i \varphi_i + f .
\]
\end{proposition}

\begin{proof}
В одну сторону переход уже проделан выше: если $u$ решает~\eqref{eq:deg-main},
то числа $u_i = (u,\psi_i)$ удовлетворяют~\eqref{eq:deg-sys}.

\medskip
Обратно, пусть $\vec u$ решает~\eqref{eq:deg-sys}. Положим
$u = \sum_j u_j \varphi_j + f$. Тогда
\[
  (u, \psi_i) = \sum_{j} u_j (\varphi_j, \psi_i) + (f,\psi_i)
              = \sum_{j} a_{ij} u_j + f_i = u_i ,
\]
где последнее равенство — это в точности~\eqref{eq:deg-sys}. Поэтому
\[
  Au = \sum_{i} (u,\psi_i)\,\varphi_i = \sum_{i} u_i \varphi_i = u - f ,
\]
то есть $u$ решает~\eqref{eq:deg-main}. Построенные отображения взаимно
обратны по построению.
\end{proof}

\subsubsection*{Сопряжённое уравнение}

Наряду с~\eqref{eq:deg-main} рассматривается сопряжённое уравнение

\stepcounter{equation}
\begin{equation}\tag{I$^*$}\label{eq:deg-adj}
  v - A^* v = g .
\end{equation}
Найдём $A^*$ явно:
\[
  (Au, v)
  = \left(\sum_{i=1}^{m} (u,\psi_i)\,\varphi_i,\ v\right)
  = \sum_{i=1}^{m} (u,\psi_i)\,(\varphi_i, v)
  = \left(u,\ \sum_{i=1}^{m} (v, \varphi_i)\,\psi_i \right) ,
\]
поскольку скаляр выносится из второго аргумента с сопряжением:
$(u,\psi_i)(\varphi_i,v) = \big(u,\ \overline{(\varphi_i,v)}\,\psi_i\big)$ и
$\overline{(\varphi_i,v)} = (v,\varphi_i)$. Итак,

\stepcounter{equation}
\begin{equation}\tag{3}\label{eq:deg-Astar}
  A^* v = \sum_{i=1}^{m} (v, \varphi_i)\, \psi_i .
\end{equation}

\begin{idea}
Обратите внимание на симметрию: у $A$ роли $\varphi$ и $\psi$ одни, у $A^*$ —
ровно обратные. Поэтому сопряжённое уравнение — это уравнение того же типа с
вырожденным ядром $\overline{K(y,x)}$, и к нему применимо всё сказанное выше.
\end{idea}

Повторяя выкладку, подставим~\eqref{eq:deg-Astar} в~\eqref{eq:deg-adj}:
\[
  v = \sum_{i=1}^{m} (v, \varphi_i)\,\psi_i + g ,
  \qquad v_i := (v, \varphi_i) ,
\]
\[
  v_i = \left(\sum_{j=1}^{m} v_j \psi_j + g,\ \varphi_i\right)
      = \sum_{j=1}^{m} v_j\, (\psi_j, \varphi_i) + (g, \varphi_i)
      = \sum_{j=1}^{m} \overline{a_{ji}}\, v_j + g_i ,
\]
так как $(\psi_j, \varphi_i) = \overline{(\varphi_i,\psi_j)}
= \overline{a_{ji}}$. В матричной записи

\stepcounter{equation}
\begin{equation}\tag{I$^{*\prime}$}\label{eq:deg-adj-sys}
  \vec v - \mathfrak{A}^*\, \vec v = \vec g ,
  \qquad \mathfrak{A}^* = \overline{\mathfrak{A}}^{\,\mathsf{T}} ,
\end{equation}
и соответствие между решениями~\eqref{eq:deg-adj} и~\eqref{eq:deg-adj-sys}
такое же, как в утверждении~\ref{prop:deg-correspondence}.

\subsubsection*{Однородные уравнения и условие ортогональности}

Полагая $f = 0$ и $g = 0$, получаем
\[
  u - Au = 0 \ \ (\mathrm{I}_0)
  \quad\Longleftrightarrow\quad
  \vec u - \mathfrak{A}\vec u = \vec 0 \ \ (\mathrm{I}_0') ,
\]
\[
  v - A^* v = 0 \ \ (\mathrm{I}_0^*)
  \quad\Longleftrightarrow\quad
  \vec v - \mathfrak{A}^*\vec v = \vec 0 \ \ (\mathrm{I}_0^{*\prime}) .
\]

Остаётся согласовать условия ортогональности в $H$ и в $\CC^m$.

\begin{lemma}\label{lem:deg-ortho}
Пусть $v$ — решение однородного сопряжённого уравнения
$(\mathrm{I}_0^*)$, а $\vec v$ — отвечающий ему вектор. Тогда
\[
  (f, v)_H = (\vec f, \vec v)_{\CC^m} .
\]
\end{lemma}

\begin{proof}
Так как $v$ решает $(\mathrm{I}_0^*)$, то по~\eqref{eq:deg-Astar}
$v = \sum_i (v,\varphi_i)\psi_i = \sum_i v_i \psi_i$. Поэтому
\[
  (f, v)_H
  = \left(f,\ \sum_{i=1}^{m} v_i \psi_i\right)
  = \sum_{i=1}^{m} \overline{v_i}\,(f, \psi_i)
  = \sum_{i=1}^{m} f_i\, \overline{v_i}
  = (\vec f, \vec v)_{\CC^m} .
\]
\end{proof}

\subsubsection*{Альтернатива для вырожденного ядра}

Теперь всё готово, чтобы перенести конечномерную
теорему~\ref{thm:fin-alternative} на интегральное уравнение.

\begin{theorem}\label{thm:deg-alternative}
Пусть ядро уравнения~\eqref{eq:deg-main} вырождено. Тогда:
\begin{enumerate}[label=\arabic*), leftmargin=2em]
  \item уравнение~\eqref{eq:deg-main} разрешимо тогда и только тогда, когда
        $f$ ортогонален всем решениям сопряжённого однородного уравнения
        $(\mathrm{I}_0^*)$, то есть $f \in \big[\ker(I - A^*)\big]^{\perp}$;
  \item \emph{(альтернатива)} либо уравнение~\eqref{eq:deg-main} однозначно
        разрешимо при любой правой части — и тогда оператор $(I-A)^{-1}$
        определён на всём $H$ и \textbf{ограничен}, — либо однородное
        уравнение $(\mathrm{I}_0)$ имеет нетривиальные решения;
  \item уравнения $(\mathrm{I}_0)$ и $(\mathrm{I}_0^*)$ имеют одинаковое
        \emph{конечное} число линейно независимых решений:
        \[
          \dim \ker (I - A) = \dim \ker (I - A^*) < \infty .
        \]
\end{enumerate}
\end{theorem}

\begin{proof}[Доказательство пунктов 1 и 2]
По утверждению~\ref{prop:deg-correspondence} решения~\eqref{eq:deg-main} и
решения~\eqref{eq:deg-sys} соответствуют друг другу, и то же верно для
сопряжённой пары. По лемме~\ref{lem:deg-ortho} условие $f \perp v$ в $H$
равносильно условию $\vec f \perp \vec v$ в $\CC^m$. Остаётся применить
утверждение~\ref{prop:fin-solvable} и теорему~\ref{thm:fin-alternative} к
системе~\eqref{eq:deg-sys}.
\end{proof}

\begin{proof}[Ограниченность $(I-A)^{-1}$]
Пусть $\det(E - \mathfrak{A}) \ne 0$. Тогда
$\vec u = (E - \mathfrak{A})^{-1}\vec f$, откуда
\[
  \abs{\vec u} \le C_1 \abs{\vec f},
  \qquad C_1 = \norm{(E-\mathfrak{A})^{-1}} .
\]
Далее, по неравенству Коши—Буняковского $\abs{f_i} = \abs{(f,\psi_i)} \le
\norm{\psi_i}\,\norm{f}$, поэтому
\[
  \abs{\vec f} \le C_2 \norm{f},
  \qquad C_2 = \Big(\sum_{i=1}^{m}\norm{\psi_i}^2\Big)^{1/2} .
\]
Наконец, из $u = \sum_i u_i \varphi_i + f$
\[
  \norm{u} \le \sum_{i=1}^{m} \norm{\varphi_i}\,\abs{u_i} + \norm{f}
  \le C_3 \abs{\vec u} + \norm{f} ,
  \qquad C_3 = \Big(\sum_{i=1}^{m}\norm{\varphi_i}^2\Big)^{1/2} ,
\]
и, собирая всё вместе,
\[
  \norm{(I-A)^{-1} f} = \norm{u} \le \big(C_1 C_2 C_3 + 1\big)\,\norm{f} .
\]
Константа зависит от систем $\{\varphi_i\}, \{\psi_i\}$ и от матрицы
$\mathfrak{A}$, но \textbf{не зависит от $f$}. Значит $(I-A)^{-1}$ ограничен.
\end{proof}

\begin{proof}[Доказательство пункта 3]
Надо проверить, что соответствие $\vec v \mapsto v = \sum_i v_i \psi_i$
сохраняет линейную независимость.

\medskip
Пусть $\vec v^{\,(1)}, \dots, \vec v^{\,(m-r)}$ — базис
$\ker(E - \mathfrak{A}^*)$, а $v^{(j)} = \sum_{i=1}^{m} v_i^{(j)} \psi_i$ —
отвечающие им решения $(\mathrm{I}_0^*)$. Пусть
\[
  \sum_{j=1}^{m-r} c_j\, v^{(j)} = 0 .
\]
Тогда
\[
  0 = \sum_{j=1}^{m-r} c_j \sum_{i=1}^{m} v_i^{(j)} \psi_i
    = \sum_{i=1}^{m} \Bigg( \sum_{j=1}^{m-r} c_j\, v_i^{(j)} \Bigg) \psi_i ,
\]
и, поскольку система $\{\psi_i\}$ линейно независима, все коэффициенты
обращаются в нуль:
\[
  \sum_{j=1}^{m-r} c_j\, v_i^{(j)} = 0, \quad i = \overline{1,m}
  \qquad\Longleftrightarrow\qquad
  \sum_{j=1}^{m-r} c_j\, \vec v^{\,(j)} = \vec 0 ,
\]
откуда $c_j = 0$ по линейной независимости $\vec v^{\,(j)}$. Значит
$v^{(1)},\dots,v^{(m-r)}$ линейно независимы, и
$\dim \ker(I - A^*) = \dim\ker(E - \mathfrak{A}^*) = m - r$.

\medskip
Точно так же — с системой $\{\varphi_i\}$ вместо $\{\psi_i\}$ — получается
\[
  \dim \ker(I - A) = \dim \ker(E - \mathfrak{A}) = m - r .
\]
Остаётся сослаться на утверждение~\ref{prop:fin-dims}.
\end{proof}

\begin{idea}
Конечность размерности здесь — не отдельное усилие, а следствие вырожденности:
все решения однородного уравнения лежат в линейной оболочке $m$ функций, так
что их не может быть больше $m$ линейно независимых. Для произвольного
компактного оператора это уже содержательное утверждение — им занимается
\S6.
\end{idea}
